\subsection{All rational factors of the pointed Fermat-eight Jacobian}

\begin{statement}{A.182}{Proved}
{Twelve positive-rank factors and an exact restriction on abelian targets}
Let \(C_8\) be the smooth proper curve
\[
              3u^8-2v^8=1,\qquad P_0=(1,1),
                    \qquad J_8=\operatorname{Jac}(C_8).
\tag{A.182.1}
\]
For nonzero rational \(A,B\), use the following notation:
\[
\begin{aligned}
 E_c&:\quad Y^2=X^3+cX,\\
 G_{A,B}&:\quad Y^2=AX^5+BX,\\
 D_{A,B}&:\quad w^4=Az^3+Bz,\qquad
 \mathcal P_{A,B}=\operatorname{im}(1-\iota)
             \subset\operatorname{Jac}(D_{A,B}),
             \quad \iota(z,w)=(z,-w).
\end{aligned}
\tag{A.182.2}
\]
The curve notation denotes the smooth proper model.  There is an
actual \(\Q\)-isogeny, constructed below from divisor pushforwards,
\[
\begin{aligned}
 J_8\longrightarrow{}& E_{-6}\times E_{18}\times E_{-12}\\
 &{}\times\prod_{(A,B)\in\mathcal G}\operatorname{Jac}(G_{A,B})
       \times\prod_{(A,B)\in\mathcal D}\mathcal P_{A,B},\\
 \mathcal G={}&\{(3,-2),(6,3),(6,-2),(24,1),(3,1),(-8,9)\},\\
 \mathcal D={}&\{(-12,1),(18,1),(3,-2)\}.
\end{aligned}
\tag{A.182.3}
\]
Every one of these twelve factors is \(\Q\)-simple and has a
rational point of infinite order.  In particular
\[
                      \operatorname{rank}J_8(\Q)\geq12.
\tag{A.182.4}
\]
Every nonzero abelian quotient of \(J_8\) over \(\Q\) has positive
rank.  More generally, for every number field \(L\) and every
abelian variety \(A/L\) with \(A(L)\) finite, every morphism from
\((C_8)_L\), or from a nonempty open subcurve of it, to \(A\) is
constant.  The last assertion is proved by an explicit Weil
restriction and Jacobian map.  No rank upper bound or classification
of sections is asserted here.
\end{statement}

\begin{proof}
\emph{The three elliptic morphisms.}
First take \(U=u^2,V=v^2\), so that
\(3U^4-2V^4=1\).  The following three pairs satisfy
\(y^2=ax^4+b\):
\[
\begin{array}{c|c|c|c}
(a,b)&x&y&(x,y)(P_0)\\ \hline
(3,-2)&u^2/v^2&1/v^4&(1,1)\\
(6,3)&v^2&3u^4&(1,3)\\
(6,-2)&u^2&2v^4&(1,2).
\end{array}
\tag{A.182.5}
\]
For each row the formula
\[
             (x,y)\longmapsto(X,Y)=(ax^2,axy)
                         \quad\hbox{in }E_{ab}
\tag{A.182.6}
\]
is an actual degree-two morphism of the proper genus-one curves.
Indeed the target equation follows by multiplying the source
equation by \(a^2x^2\); its function field determines \(x^2\) and
then \(y=Y/(ax)\), while \((x,y)\mapsto(-x,-y)\) is the
nonidentity automorphism fixing both target functions.  Translation
by the negative of the displayed image point gives the pointed
isogeny when an origin on the quartic is wanted.  The morphisms
from \(C_8\) and their divisor pushforwards require no such choice.

\emph{The six genus-two morphisms.}
The intermediate functions in the next table satisfy
\(y^2=Ax^8+B\).  Composing them with
\((X,Y)=(x^2,xy)\) gives the stated morphism to \(G_{A,B}\):
\[
\begin{array}{c|c|c|c}
(A,B)&x&y&(X,Y)(P_0)\\ \hline
(3,-2)&u/v&1/v^4&(1,1)\\
(6,3)&v&3u^4&(1,3)\\
(6,-2)&u&2v^4&(1,2)\\
(24,1)&uv&3u^8+2v^8&(1,5)\\
(3,1)&u/v^2&(3u^8+1)/(2v^8)&(1,2)\\
(-8,9)&v/u^2&(2v^8-1)/u^8&(1,1).
\end{array}
\tag{A.182.7}
\]
Substitution of \(3u^8-2v^8=1\) verifies every row.  For example,
\((3u^8+2v^8)^2=24u^8v^8+1\), and
\((3u^8+1)^2-12u^8=4v^{16}\).  The final composition satisfies
\(Y^2=AX^5+BX\) literally.  All these rational maps are
nonconstant and extend uniquely to the proper curves.

\emph{The remaining three surface morphisms.}
There are actual morphisms to the following plane quartics:
\[
\begin{array}{c|c|c|c}
(A,B)&z&w&(z,w)(P_0)\\ \hline
(-12,1)&-u^4/2&uv^2&(-1/2,1)\\
(18,1)&v^4/3&u^2v&(1/3,1)\\
(3,-2)&u^4/v^4&u/v^3&(1,1).
\end{array}
\tag{A.182.8}
\]
Each row satisfies \(w^4=Az^3+Bz\).  In projective coordinates
the target is
\[
                      W^4=AZ^3T+BZT^3.
\tag{A.182.9}
\]
It is smooth in characteristic zero and at every prime not
dividing \(2AB\).  At a point with \(W\ne0\) use the
\(W\)-derivative.  At \(W=0\), the binary quartic on the right
has four distinct roots.  Its genus is three, and its unique point
at infinity is rational.  The involution \(\iota\) fixes that
point.  The quotient map
\[
       \pi:D_{A,B}\longrightarrow E'_{A,B},\qquad
       (z,w)\longmapsto(z,w^2),\qquad
       E'_{A,B}:y^2=Az^3+Bz
\tag{A.182.10}
\]
has degree two.  On Jacobians,
\(\pi_*\pi^*=[2]\) and
\(\pi^*\pi_*=1+\iota\).  Consequently
\(\mathcal P_{A,B}=\ker(\pi_*)^0\) has dimension two, and
the actual homomorphism
\[
 \operatorname{Jac}(D_{A,B})\longrightarrow
       E'_{A,B}\times\mathcal P_{A,B},\qquad
                 x\longmapsto(\pi_*x,(1-\iota)x)
\tag{A.182.11}
\]
is an isogeny.  Its kernel is killed by two: if both entries
vanish, then \(\iota x=x\) and
\(2x=(1+\iota)x=\pi^*\pi_*x=0\).
Composition of each divisor pushforward in (A.182.8) with
\(1-\iota\) defines the three surface components of (A.182.3).

\emph{The full differential calculation.}
The smooth plane curve (A.182.1) has genus 21.  A basis of regular
differentials is
\[
 \omega_{a,b}=u^{a-1}v^{b-8}\,du,
          \qquad a,b\geq1,\quad a+b\leq7.
\tag{A.182.12}
\]
This is the plane-curve differential basis obtained by dividing
monomials of degree at most five by the \(v\)-derivative of the
defining equation; the omitted common factor is \(-16\).
Under \((u,v)\mapsto(\zeta u,\xi v)\), its character is
\(\zeta^a\xi^b\).

For each genus-two target use \(dX/Y,X\,dX/Y\).  For a quartic
\(D_{A,B}\), its three regular differentials are
\(dz/w^2,dz/w^3,z\,dz/w^3\).  Their orders at infinity are
\(1,4,0\); finite ramification points also give nonnegative
orders.  The last two are anti-invariant under \(\iota\).
If \(j:\mathcal P_{A,B}\hookrightarrow\operatorname{Jac}(D_{A,B})\)
is inclusion, put \(\eta_0=\tfrac12j^*(dz/w^3)\) and
\(\eta_1=\tfrac12j^*(z\,dz/w^3)\).  The pullbacks under
\(1-\iota\) are exactly the two displayed anti-invariant forms.
Thus the following table records the literal cotangent maps of
(A.182.3).  A row \((a,b);c\) means pullback
\(c\omega_{a,b}\):
\[
\begin{array}{c|c|c}
\text{target}&\text{first differential}&\text{second differential}\\ \hline
E_{-6}&(2,2);-1&\\
E_{18}&(4,2);1&\\
E_{-12}&(2,4);1&\\ \hline
G_{3,-2}&(1,3);-1&(3,1);-1\\
G_{6,3}&(4,1);1&(4,3);1\\
G_{6,-2}&(1,4);1&(3,4);1\\
G_{24,1}&(1,1);1&(3,3);1\\
G_{3,1}&(1,6);-2&(3,2);-2\\
G_{-8,9}&(6,1);-1&(2,3);-1\\ \hline
\mathcal P_{-12,1}&(1,2);-2&(5,2);1\\
\mathcal P_{18,1}&(2,1);2&(2,5);2/3\\
\mathcal P_{3,-2}&(1,5);-2&(5,1);-2.
\end{array}
\tag{A.182.13}
\]
The elliptic differential is \(dX/(2Y)\).  Every entry follows
by differentiating the literal functions above with
\(dv/du=3u^7/(2v^7)\); the retained certificate checks the
resulting rational-function identity modulo (A.182.1), including
the scalar.  The table contains each of the 21 indices in
(A.182.12) once.  Therefore the product of the actual
pushforward homomorphisms has invertible cotangent map.
Its kernel has dimension zero and its image has dimension 21;
properness then makes it a surjective isogeny.  This proves
(A.182.3), including exhaustion of all rational factors.

\emph{Exact good-reduction counts.}
For a genus-two curve let \(N_i\) be its number of points over
\(\mathbf F_{p^i}\), and put
\[
 t_1=p+1-N_1,\quad t_2=p^2+1-N_2,\qquad
 P_p(T)=T^4-t_1T^3+\frac{t_1^2-t_2}{2}T^2-pt_1T+p^2.
\tag{A.182.14}
\]
Then \(\#J(\mathbf F_p)=P_p(1)\).  These are the curve trace
and abelian-variety point formulas in
\cite[III, Theorem 11.1; II, Theorem 1.1]{LegMilneAV2008}.
For a Prym in (A.182.11), subtract the elliptic traces from the
quartic traces before using (A.182.14); this is the Tate-module
map of that actual isogeny.

The local certificate enumerates the fields themselves.  For
\(\mathbf F_{p^2}=\mathbf F_p[d]/(d^2-\nu)\), where \(\nu\)
is the least quadratic nonresidue, the quadratic character of
\(a+bd\) is the Legendre symbol of \(a^2-\nu b^2\).
A nonzero fourth-power equation has \(\gcd(4,q-1)\) roots
exactly when its right side raised to
\((q-1)/\gcd(4,q-1)\) is one.  Each model here has one point
at infinity.  These formulas give, in particular,
\[
\begin{array}{c|c|c|c|c}
\text{surface}&p&N_1&N_2&P_p(T)\\ \hline
J(G_{3,-2})&13&14&170&T^4+169\\
J(G_{6,3})&5&6&26&T^4+25\\
J(G_{6,-2})&5&6&26&T^4+25\\
J(G_{24,1})&13&14&170&T^4+169\\
J(G_{3,1})&5&6&26&T^4+25\\
J(G_{-8,9})&5&6&26&T^4+25.
\end{array}
\tag{A.182.15}
\]
For the Pryms, the two point-count pairs below are respectively
those of \(D_{A,B}\) and its elliptic quotient:
\[
\begin{array}{c|c|c|c|c}
(A,B)&p&(N_1(D),N_2(D))&(N_1(E'),N_2(E'))&P_p(T)\\ \hline
(-12,1)&5&(10,36)&(10,20)&T^4+8T^2+25\\
(18,1)&5&(10,36)&(10,20)&T^4+8T^2+25\\
(3,-2)&13&(18,132)&(18,180)&T^4-24T^2+169.
\end{array}
\tag{A.182.16}
\]
All primes used are good.  For the hyperelliptic models this
follows from the nonzero discriminant of \(AX^5+BX\), whose
prime support is contained in \(\{2,3\}\).  The quartic
smoothness check was given after (A.182.9).  At these odd primes
the involution and the degree-two quotient extend over the
smooth models.  On their good Jacobians the connected kernel
of \(\pi_*\) is a smooth proper group scheme, because
\(\pi_*\pi^*=[2]\) is etale.  It is the good Prym model.
Thus all reductions in (A.182.15)--(A.182.16) are the reductions
of the stated abelian varieties.

\emph{Torsion bounds and infinite-order points.}
The following pairs of good group orders suffice:
\[
\begin{array}{c|c|c}
\text{factor}&(p,\#A(\mathbf F_p))&(p',\#A(\mathbf F_{p'}))\\ \hline
E_{-6}&(7,8)&(13,18)\\
E_{18},E_{-12}&(5,10)&(7,8)\\
J(G_{3,-2}),J(G_{24,1})&(7,64)&(13,170)\\
J(G_{6,3}),J(G_{6,-2}),J(G_{3,1}),J(G_{-8,9})&(5,26)&(7,64)\\
\mathcal P_{-12,1},\mathcal P_{18,1}&(5,34)&(7,64)\\
\mathcal P_{3,-2}&(5,32)&(13,146).
\end{array}
\tag{A.182.17}
\]
Reduction injects prime-to-residue-characteristic torsion:
the formal reduction kernel is pro-\(p\), with multiplication
by every integer prime to \(p\) invertible.  Each row therefore
bounds rational torsion by order two.  At either residue
characteristic, use the other column; that column has order
prime to that characteristic.  This accounts for every prime.

On the elliptic factors the points
\[
                  (3,3)\in E_{-6},\qquad
                  (6,18)\in E_{18},\qquad
                  (6,12)\in E_{-12}
\tag{A.182.18}
\]
have nonzero second coordinate, so are not two-torsion and
have infinite order.  On each genus-two factor use
\(R=[(1,y_0)-\infty]\), with \(y_0\) from (A.182.7).
The space \(L(2\infty)\) is exactly
\(\langle1,X\rangle\): \(X\) has pole order two and
\(Y\) pole order five, and parity prevents cancellation
between the polynomial and \(Y\)-polynomial terms.
No such function has divisor \(2(1,y_0)-2\infty\), since
\(y_0\ne0\) and a zero of a linear polynomial in \(X\)
also vanishes at \((1,-y_0)\).  Hence \(2R\ne0\), so
\(R\) has infinite order.

For each Prym let \(P=(z_0,1)\) be the rational point in
(A.182.8).  The class
\[
       R=[P-\iota P]=(1-\iota)[P-\infty]
                          \quad\hbox{belongs to }\mathcal P_{A,B}(\Q).
\tag{A.182.19}
\]
The canonical map of (A.182.9), in the three differential
coordinates used above, is \([w:1:z]\), its plane-quartic
embedding.  A degree-two rational function would have a
nonidentity involution whose invariant regular differentials
descend to \(\mathbf P^1\) and vanish.  It would therefore
identify its conjugate points under this canonical map,
contrary to that embedding.  Thus no degree-two function
exists.  Since \(P\ne\iota P\), a relation
\(2[P-\iota P]=0\) would produce precisely such a function.
This proves \(2R\ne0\), and the order bound in (A.182.17)
again proves infinite order.  The retained certificate also
checks this nonvanishing by a literal full-rank jet matrix.

\emph{Simplicity and the rank conclusion.}
All quartics in (A.182.15)--(A.182.16) are irreducible over
\(\Q\).  For \(T^4+aT^2+p^2\), a factorization into
monic quadratics has opposite linear coefficients.  If these
coefficients vanish, \(a^2-4p^2\) is a rational square.
Otherwise the two constant terms coincide, equal \(p\) or
\(-p\), forcing respectively \(2p-a\) or \(-2p-a\)
to be a rational square.  A rational linear root also gives
a quadratic factor by adjoining its negative.  For the
four possibilities used here, the three necessary squares are
\[
\begin{array}{c|c}
(a,p)&a^2-4p^2,\ 2p-a,\ -2p-a\\ \hline
(0,5)&-100,10,-10\\
(0,13)&-676,26,-26\\
(8,5)&-36,2,-18\\
(-24,13)&-100,50,-2.
\end{array}
\tag{A.182.20}
\]
None is a rational square.  A proper nonzero rational abelian
subvariety of any of the nine surfaces and its quotient would
have good reduction at the chosen prime and would factor the
actual Frobenius polynomial into integral positive-degree
polynomials.  This follows from the exact rational Tate-module
sequence and the good-reduction criterion
\cite[IV, Theorem 3.5 and Corollary 3.6]{LegMilneAV2008}.
It contradicts (A.182.20).  The nine surfaces are therefore
\(\Q\)-simple; the elliptic factors are simple by dimension.

An isogeny admits an inverse up to multiplication by a nonzero
integer, by factoring that multiplication through its finite
kernel.  It therefore induces an isomorphism after tensoring
rational point groups with \(\Q\).  The twelve explicit
infinite-order points give (A.182.4).  Finite generation, when
expressing this in terms of rank, is
\cite[I, Theorem 16.7]{LegMilneAV2008}.

\emph{All number-field targets with finitely many rational points.}
Let \(B\) be a nonzero rational abelian quotient of \(J_8\).
Choose an inverse to (A.182.3) up to multiplication.  Its
composition with \(J_8\to B\) is nonzero on at least one
of the twelve simple factors.  On that factor its kernel has
dimension zero, by simplicity.  The explicit infinite-order
point of that factor therefore maps to an infinite-order
point of \(B(\Q)\).  This proves the quotient assertion
without assuming that the twelve factors are pairwise
nonisogenous.

Now suppose \(L/\Q\) is finite and \(f:(C_8)_L\to A/L\)
is nonconstant.  Replace \(f\) by \(f-f(P_0)\).  For every
\(\Q\)-scheme \(S\), the functorial map
\[
 C_8(S)\longrightarrow C_8(S\times_{\Q}L)
       \xrightarrow{\ f\ }A(S\times_{\Q}L)
       =\bigl(\operatorname{Res}_{L/\Q}A\bigr)(S)
\tag{A.182.21}
\]
gives the actual adjoint morphism \(\widetilde f\) over
\(\Q\).  After base change to an algebraic closure its
components are the conjugates of \(f\), so it is nonconstant.
It sends \(P_0\) to zero.  The Jacobian universal property
\cite[III, Proposition 6.1]{LegMilneAV2008} supplies a nonzero
homomorphism
\[
         J_8\longrightarrow\operatorname{Res}_{L/\Q}A.
\tag{A.182.22}
\]
Its image \(B\) is a nonzero rational abelian quotient of
\(J_8\).  We just proved that \(B(\Q)\) contains a point
of infinite order, whereas
\[
                   B(\Q)\subset
              \bigl(\operatorname{Res}_{L/\Q}A\bigr)(\Q)=A(L).
\tag{A.182.23}
\]
This excludes finite \(A(L)\).  A morphism from a nonempty
open subcurve extends across each missing point by the
valuative criterion for the proper target \(A\), so the
same proof applies to the open curve.

The fully local, executed certificate is
\path{code/point3_fermat8_decomposition_certificate.py}, with
retained output and data.  It re-executes the two companion
surface certificates, checks all literal maps and 21 exact
differential scalars, all torsion bounds, and every stated
irreducibility assertion.  It uses no external rank calculation.
The exact ranks are determined in A.185--A.186 below. The new
factors at level 24 remain to be determined by further work.
\end{proof}
